% ═══════════════════════════════════════════════════════════════
%  ch02.tex — فصل ۲
%  نجات‌بخشی در برابر نجات‌طلبی: تحلیل چندسویه
%  نسخه‌ی اصلاح‌شده: رفع مشکلات RTL، رنگ و جهت
% ═══════════════════════════════════════════════════════════════

\chapter{نجات‌بخشی در برابر نجات‌طلبی: تحلیل چندسویه}
\label{ch:rescue-vs-seeking}

\begin{tcolorbox}[quotebox, title={}]
	«آزادی‌ای که به مردم اهدا شود، از آزادی حقیقی چیزی ندارد.
	باید آزادی را خود به چنگ آورد، نه آنکه دیگری آن را تقدیمت
	کند؛ زیرا آنچه تقدیم می‌شود، به همان آسانی پس گرفته
	می‌شود.»
	
	\hfill ---فریدریش شیلر\src{\lr{Schiller, \textit{Wilhelm Tell}, 1804}}
\end{tcolorbox}


% ────────────────────────────────────────────────────────────────
\section{درآمد: چرا این فصل ضروری است؟}
\label{sec:ch02-intro}

فصل پیشین مفاهیم پایه را تبیین کرد. اما پیچیده‌ترین بخش بحث
مداخله‌ی خارجی نه در تعریف مفاهیم، بلکه در \concept{تنش‌های درونی و
	پارادوکس‌هایی} است که هنگام تحقق آن‌ها ظاهر می‌شوند. این فصل به
هشت تنش بنیادین می‌پردازد که هر یک سویه‌ای متفاوت از مسئله‌ی
نجات‌بخشی و نجات‌طلبی را آشکار می‌کنند.

\begin{tcolorbox}[mybox, title={ساختار فصل: هشت تنش بنیادین}]
	\begin{enumerate}[label=\textbf{تنش \arabic*:}, nosep, rightmargin=1em]
		\item امکان رهایی‌بخشی ساختاری از بیرون
		\item پارادوکس اجبار به آزادی
		\item تنش میان مدرنیته‌ی خطی و تکثر فرهنگی
		\item جدایی جبرانی و زنجیره‌ی خشونت
		\item گروه پیشرو، بازیگر خارجی و مشروعیت مداخله
		\item بلوغ درون‌زاد ارزش‌ها و ناممکنی تزریق
		\item جامعه‌ی دوپاره و شکنندگی دموکراسی
		\item بازی‌ای که نمی‌توان متوقف کرد: «گر رسم شود که مست گیرند»
	\end{enumerate}
\end{tcolorbox}


% ════════════════════════════════════════════════════════════════
\section{تنش اول: امکان رهایی‌بخشی ساختاری از بیرون}
\label{sec:ch02-tension1}

\subsection{طرح مسئله}
\label{subsec:ch02-t1-problem}

آیا اساساً ممکن است که یک نیروی خارجی، به‌خواست یا بی‌خواست مردم،
تغییراتی ساختاری ایجاد کند که واقعاً به \concept{کاهش رنج}،
\concept{رفع سلطه} و \concept{توانمندسازی} منجر شود؟ پاسخ تجربی
آن است: بله، اما مشروط.

\subsection{شواهد مثبت: مداخله‌ی ساختارساز}
\label{subsec:ch02-t1-positive}

مهم‌ترین نمونه‌ی تاریخی، \concept{بازسازی ژاپن و آلمان} پس از
جنگ جهانی دوم است. ژنرال
\thinker{داگلاس مک‌آرتور}%
\footnote{%
	داگلاس مک‌آرتور (\lr{Douglas MacArthur, 1880--1964}).%
}%
در مقام فرماندار نظامی ژاپن (۱۹۴۵--۱۹۵۱) اصلاحاتی ساختاری انجام
داد که عبارت بودند از:

\begin{enumerate}[label=\textbf{\alph*.}, nosep, rightmargin=1em]
	\item تدوین قانون اساسی جدید با تضمین حقوق فردی و دموکراسی
	پارلمانی\src{\lr{Dower, 1999, pp.\,346--373}}؛
	\item اصلاحات ارضی گسترده که ساختار فئودالی را درهم
	شکست\src{\lr{ibid., pp.\,525--546}}؛
	\item بهداشت عمومی و آموزش اجباری که به بهبود چشمگیر
	شاخص‌های انسانی انجامید؛
	\item تفکیک دین از دولت (حذف شینتوی دولتی).
\end{enumerate}

\begin{tcolorbox}[casebox,
	title={مطالعه‌ی موردی: بهداشت به‌مثابه‌ی رهایی‌بخشی}]
	
	آیا \concept{واکسیناسیون اجباری} شکلی از اجبار است؟ بی‌تردید.
	اما آیا رهایی‌بخش نیز هست؟ تاریخ پاسخ مثبت می‌دهد. برنامه‌ی
	ریشه‌کنی آبله توسط سازمان بهداشت جهانی (۱۹۶۷--۱۹۸۰) نمونه‌ی
	بارزی است از «اجباری که در دامان خود آزادی
	دارد»\src{\lr{Fenner et al., 1988, pp.\,421--538}}.
	
	امنیت نیز چنین است: ممانعت از ورود شهروندان به مناطق مین‌گذاری‌شده
	یا بمباران‌شده «اجباری رهایی‌بخش» است. در این موارد، \concept{اجبار}
	نه در تعارض با \concept{خودمختاری}%
	\footnote{خودمختاری؛ به انگلیسی: \lr{Autonomy}}%
	بلکه مقوّم آن است؛ زیرا پیش‌شرط خودمختاری، \concept{زنده ماندن
		و سلامت} است\src{\lr{Sen, 1999, pp.\,15--30}}.
\end{tcolorbox}

\subsection{مرزهای امکان: چه وقت ساختارسازی بیرونی کار می‌کند؟}
\label{subsec:ch02-t1-conditions}

\thinker{آمارتیا سن}%
\footnote{%
	آمارتیا سن (\lr{Amartya Sen, 1933–})؛%
	اقتصاددان و فیلسوف هندی، برنده‌ی نوبل اقتصاد ۱۹۹۸.%
}%
در نظریه‌ی \concept{قابلیت‌ها}%
\footnote{رویکرد قابلیت‌ها؛ به انگلیسی: \lr{Capabilities Approach}}%
استدلال می‌کند که «آزادی مثبت»، یعنی توانایی واقعی انسان برای
زیستن به شیوه‌ای که ارزشمند می‌داند، مستلزم پیش‌شرط‌های مادی
و نهادی است\src{\lr{Sen, 1999, pp.\,13--34}}. این پیش‌شرط‌ها
(بهداشت، آموزش، امنیت، تغذیه) ممکن است از بیرون فراهم شوند،
مشروط بر آنکه:

\begin{enumerate}[label=\textbf{\arabic*.}, nosep, rightmargin=1em]
	\item ساختارسازی با مشارکت محلی و نه صرفاً تحمیلی باشد؛
	\item نیروی خارجی برنامه‌ی خروج مشخص و تعهد به انتقال
	اختیارات داشته باشد؛
	\item اهداف ساختاری بر منافع استراتژیک مداخله‌گر غلبه
	داشته باشد یا حداقل با آن‌ها همسو باشد؛
	\item زمان کافی برای نهادینه‌شدن تغییرات در نظر گرفته شود.
\end{enumerate}

\begin{tcolorbox}[enemybox, title={نمونه‌ی ناموفق: عراق ۲۰۰۳}]
	مقایسه‌ی ژاپن ۱۹۴۵ با عراق ۲۰۰۳ آموزنده است. در ژاپن، اشغالگر
	(آمریکا) با بوروکراسی کارآمد محلی همکاری کرد، برنامه‌ی بلندمدت
	داشت و جامعه‌ی نسبتاً همگنی را اصلاح می‌کرد. در عراق، ساختار
	بوروکراتیک منحل شد (\lr{De-Ba'athification})، برنامه‌ی خروج
	وجود نداشت و جامعه به‌شدت چندپاره بود. نتیجه: «ساختارسازی»
	به ساختارشکنی تبدیل
	شد\src{\lr{Diamond, 2005, pp.\,280--315; Dodge, 2009, pp.\,23--68}}.
\end{tcolorbox}


% ════════════════════════════════════════════════════════════════
\section{تنش دوم: پارادوکس اجبار به آزادی}
\label{sec:ch02-tension2}

\subsection{طرح پارادوکس}
\label{subsec:ch02-t2-paradox}

\begin{tcolorbox}[defbox, title={پارادوکس اجبار به آزادی}]
	آیا می‌توان کسی را \textbf{مجبور به آزاد بودن} کرد؟ اگر آزادی
	ذاتاً مستلزم \concept{خودمختاری} و \concept{انتخاب آگاهانه} است،
	آیا فرایندی که با اجبار آغاز شود می‌تواند به آزادی واقعی ختم شود؟
\end{tcolorbox}

این پرسش ریشه‌ای فلسفی دارد.
\thinker{ژان-ژاک روسو}%
\footnote{%
	ژان-ژاک روسو (\lr{Jean-Jacques Rousseau, 1712--1778}).%
}%
در کتاب قرارداد اجتماعی عبارت بحث‌برانگیز «مجبور کردن
به آزادی» را به کار برده
است\src{\lr{Rousseau, 1762/2002, I.vii}}.
\thinker{آیزایا برلین}%
\footnote{%
	آیزایا برلین (\lr{Isaiah Berlin, 1909--1997}).%
}%
این ایده را خطرناک دانسته و هشدار داده که «اجبار به آزادی» می‌تواند
توجیه‌گر بدترین اشکال استبداد
شود\src{\lr{Berlin, 1969, pp.\,131--134}}.

\subsection{آزادی مثبت و پیش‌شرط‌های آن}
\label{subsec:ch02-t2-positive}

اما مسئله پیچیده‌تر از یک بله/خیر ساده است. تمایز برلین
میان \concept{آزادی منفی}%
\footnote{آزادی منفی (\lr{Negative Liberty}): آزادی \textit{از}
	موانع بیرونی.}%
و \concept{آزادی مثبت}%
\footnote{آزادی مثبت (\lr{Positive Liberty}): آزادی \textit{برای}
	تحقق قابلیت‌ها.}%
راه‌گشاست\src{\lr{Berlin, 1969, pp.\,118--172}}.

اجبار به آموزش، واکسیناسیون، رعایت مقررات ایمنی و حتی پرداخت
مالیات، نمونه‌هایی از \textbf{محدودیت آزادی منفی در خدمت آزادی
	مثبت} هستند. این‌ها در سطح ملی پذیرفته‌شده‌اند. پرسش آن است:
آیا همین منطق در سطح بین‌المللی نیز اعمال‌شدنی است؟

\begin{figure}[htbp]
	\centering
	\begin{tikzpicture}[
		>=Stealth,
		every node/.style={font=\small}
		]
		% دو نیم‌دایره
		\fill[BgTeal!40] (0,0) arc (180:0:3) -- cycle;
		\fill[red!10] (0,0) arc (180:360:3) -- cycle;
		
		% خطوط
		\draw[very thick, PrimaryDeep] (0,0) arc (180:0:3);
		\draw[very thick, AccentRed] (0,0) arc (180:360:3);
		\draw[thick, AccentGold, dashed] (-0.5,0) -- (6.5,0);
		
		% برچسب‌ها
		\node[text=AccentTeal, font=\bfseries] at (3,1.5)
		{\rl{آزادی مثبت}};
		\node[text=AccentTeal, font=\small] at (3,0.8)
		{\rl{توانمندسازی، آموزش، بهداشت}};
		
		\node[text=AccentRed, font=\bfseries] at (3,-1.5)
		{\rl{اجبار}};
		\node[text=AccentRed, font=\small] at (3,-0.8)
		{\rl{محدودیت، الزام، مداخله}};
		
		% فلش دوطرفه
		\draw[<->, very thick, AccentGold] (6.8,-2) -- (6.8,2)
		node[midway, right, font=\small\bfseries, text=AccentAmber,
		text width=2.5cm, align=center]
		{\rl{مرز مشروعیت}\\[2pt]{\scriptsize\rl{کجاست؟}}};
		
		% عنوان
		\node[above, font=\normalsize\bfseries, text=PrimaryDeep]
		at (3,3.5) {\rl{تنش میان اجبار و آزادی مثبت}};
	\end{tikzpicture}
	\caption{نمودار مفهومی تنش میان اجبار و آزادی مثبت}
	\label{fig:ch02-coercion-freedom}
\end{figure}

\subsection{سه شرط مشروعیت اجبار رهایی‌بخش}
\label{subsec:ch02-t2-conditions}

بر اساس ترکیب آرای
\thinker{جان استوارت میل}%
\footnote{جان استوارت میل (\lr{John Stuart Mill, 1806--1873}).}%
(اصل ضرر)،
\thinker{آمارتیا سن} (رویکرد قابلیت‌ها) و
\thinker{مارتا نوسباوم}%
\footnote{%
	مارتا نوسباوم (\lr{Martha C.\,Nussbaum, 1947–}).%
}%
(فهرست قابلیت‌های
مرکزی)\src{\lr{Nussbaum, 2011, pp.\,33--34}}،
سه شرط برای مشروعیت «اجبار رهایی‌بخش» پیشنهاد می‌کنیم:

\begin{enumerate}[label=\textbf{\arabic*.}, nosep, rightmargin=1em]
	\item \concept{اصل حداقلی بودن}: اجبار باید در حداقل ممکن
	اعمال شود و فقط ناظر بر پیش‌شرط‌های بنیادین حیات و کرامت
	باشد (جان، سلامت، آموزش پایه)؛
	\item \concept{اصل موقتی بودن}: اجبار باید خصلت انتقالی داشته
	باشد و با توانمند شدن جامعه‌ی هدف، به‌تدریج کنار رود؛
	\item \concept{اصل قابلیت‌محوری}: هدف اجبار نه تحمیل محتوای
	خاص فرهنگی، بلکه فراهم‌سازی \textit{امکان انتخاب} باشد؛
	یعنی «قابلیت» به معنای سِنی کلمه.
\end{enumerate}

\begin{tcolorbox}[wavebox, title={نکته‌ی انتقادی}]
	حتی این سه شرط نیز از نقد مصون نیستند.
	\thinker{طلال اسد}%
	\footnote{%
		طلال اسد (\lr{Talal Asad, 1932--2024})؛%
		مردم‌شناس و نظریه‌پرداز مطالعات پسااستعماری.%
	}%
	استدلال می‌کند که مفهوم «قابلیت‌های بنیادین» خود بر فرض‌های
	فرهنگی خاصی (عمدتاً لیبرال-غربی) استوار است و ادعای جهان‌شمولی
	آن نوعی \concept{امپریالیسم معرفتی}%
	\footnote{امپریالیسم معرفتی؛ به انگلیسی: \lr{Epistemic Imperialism}}%
	است\src{\lr{Asad, 2003, pp.\,127--158}}.
\end{tcolorbox}


% ════════════════════════════════════════════════════════════════
\section{تنش سوم: مدرنیته‌ی خطی در برابر تکثر فرهنگی}
\label{sec:ch02-tension3}

\subsection{اجبار به پذیرش مفاهیم مدرن}
\label{subsec:ch02-t3-coercion}

هنگامی که نیروی خارجی (یا حتی نخبگان داخلی متأثر از بیرون)
می‌خواهند جامعه‌ای را «مدرن» کنند، ناگزیر مجموعه‌ای از ارزش‌ها
و نهادها را ترویج می‌کنند که ریشه در تجربه‌ی تاریخی اروپا
دارند: فردگرایی، سکولاریسم، حقوق بشر جهانی، دموکراسی نمایندگی.

\thinker{دیپش چاکراوارتی}%
\footnote{%
	دیپش چاکراوارتی (\lr{Dipesh Chakrabarty, 1948–})؛%
	مورخ هندی‌تبار و نویسنده‌ی%
	\lr{\textit{Provincializing Europe}} (۲۰۰۰).%
}%
استدلال می‌کند که این رویکرد مبتنی بر نگرش
\concept{توسعه‌ی خطی}%
\footnote{توسعه‌ی خطی؛ به انگلیسی:
	\lr{Linear Development / Stadial Theory}}%
و \concept{اروپامحوری}%
\footnote{اروپامحوری؛ به انگلیسی: \lr{Eurocentrism}}%
است: همه‌ی جوامع باید همان مسیری را طی کنند که اروپا طی کرده
است\src{\lr{Chakrabarty, 2000, pp.\,3--23}}.

\subsection{بردگان خشنود و زنان خوگرفته: مسئله‌ی ترجیحات تطبیقی}
\label{subsec:ch02-t3-adaptive}

\begin{tcolorbox}[defbox, title={مفهوم: ترجیحات تطبیقی}]
	\concept{ترجیحات تطبیقی}%
	\footnote{ترجیحات تطبیقی؛ به انگلیسی: \lr{Adaptive Preferences}}%
	به ترجیحاتی اطلاق می‌شود که افراد در پاسخ به محدودیت‌های موجود
	شکل می‌دهند؛ نه بر اساس آنچه واقعاً ممکن یا مطلوب است، بلکه
	بر اساس آنچه در دسترس‌شان قرار گرفته. سیاهانی که با تبعیض نژادی
	«کنار آمده‌اند»، زنانی که فرهنگ پدرسالار را درونی و حتی ترویج
	کرده‌اند، و بردگانی که «خوشحال» به نظر می‌رسند، نمونه‌های
	کلاسیک این
	پدیده‌اند\src{\lr{Elster, 1983, pp.\,109--140; Nussbaum, 2000, pp.\,111--166}}.
\end{tcolorbox}

\thinker{یون الستر}%
\footnote{%
	یون الستر (\lr{Jon Elster, 1940–})؛%
	فیلسوف و نظریه‌پرداز انتخاب عقلانی نروژی.%
}%
با داستان روباه و انگورهای ترش ایسوپ، نشان می‌دهد که انسان‌ها
اغلب ارزش چیزهای دست‌نیافتنی را کاهش می‌دهند تا ناهماهنگی
شناختی را کم کنند\src{\lr{Elster, 1983, pp.\,109--140}}.
پرسش اخلاقی آن است: آیا احترام به «ترجیحات» این افراد احترام به
آزادی آن‌هاست، یا مشارکت در ستم بر آن‌ها؟

\subsection{نقد چندفرهنگی‌گرایی و نقد نقد}
\label{subsec:ch02-t3-multiculturalism}

\begin{tcolorbox}[comparebox,
	title={دو خوانش رقیب از تکثر فرهنگی}]
	\begin{tabular}{>{\bfseries\color{PrimaryMid}}p{2.5cm}
			p{5cm} p{5cm}}
		\toprule
		& \textbf{جهان‌شمولی حقوق بشر} &
		\textbf{نسبی‌گرایی فرهنگی} \\
		\midrule
		ادعای اصلی &
		برخی حقوق جهان‌شمول‌اند و تابع فرهنگ نیستند &
		حقوق همواره در بافت فرهنگی معنا می‌یابند \\
		\midrule
		نماینده &
		نوسباوم، سن، اعلامیه‌ی جهانی حقوق بشر &
		چاکراوارتی، اسد، مکتب پسااستعماری \\
		\midrule
		قوت &
		مبنای مشترک برای نقد ستم در هر فرهنگ &
		احترام به تنوع و مقاومت در برابر سلطه‌ی فرهنگی \\
		\midrule
		ضعف &
		خطر تحمیل مدل غربی &
		خطر توجیه ستم به‌نام «فرهنگ» \\
		\bottomrule
	\end{tabular}
\end{tcolorbox}

راه برون‌رفت شاید در آن چیزی باشد که
\thinker{سیلا بن‌حبیب}%
\footnote{%
	سیلا بن‌حبیب (\lr{Seyla Benhabib, 1950–})؛%
	فیلسوف سیاسی ترکیه‌ای-آمریکایی.%
}%
\concept{جهان‌شمولی گفت‌وگویی}%
\footnote{جهان‌شمولی گفت‌وگویی؛ به انگلیسی:
	\lr{Dialogical Universalism}}%
می‌نامد: ارزش‌های جهان‌شمول نه از بالا تحمیل می‌شوند و نه
نسبی‌اند، بلکه از طریق گفت‌وگوی برابر میان فرهنگ‌ها استخراج
می‌شوند\src{\lr{Benhabib, 2002, pp.\,26--48}}.


% ════════════════════════════════════════════════════════════════
\section{تنش چهارم: جدایی جبرانی، حق تعیین سرنوشت و زنجیره‌ی خشونت}
\label{sec:ch02-tension4}

\subsection{فرهنگ‌های اجبارگر و حق جدایی}
\label{subsec:ch02-t4-oppressive}

در برخی جوامع، ساختار قدرت نه‌تنها سرکوبگرانه است، بلکه
\textbf{بر گفتمان اجبار، کیفر و سختگیری بنا شده}. در این‌گونه
نظام‌ها، آنچه بر همگان تحمیل می‌شود نه یک «فرهنگ» به معنای
مردم‌شناختی، بلکه یک \concept{برساخته‌ی قدرت}%
\footnote{برساخته‌ی قدرت؛ به انگلیسی: \lr{Power Construct}}%
است که ذینفعان آن از لوای «سنت»، «دین» یا «هویت ملی» برای
مشروعیت‌بخشی استفاده می‌کنند.

برای کسانی که در درون چنین نظامی هستند و آن را نمی‌پذیرند، آیا
\concept{حق تعیین سرنوشت}%
\footnote{حق تعیین سرنوشت؛ به انگلیسی:
	\lr{Right to Self-Determination}}%
شامل حق خروج و جدایی نیز هست؟

\subsection{جدایی جبرانی در حقوق بین‌الملل}
\label{subsec:ch02-t4-remedial}

\begin{tcolorbox}[legalbox, title={مفهوم: جدایی جبرانی}]
	\concept{جدایی جبرانی}%
	\footnote{جدایی جبرانی؛ به انگلیسی: \lr{Remedial Secession}}%
	نظریه‌ای در حقوق بین‌الملل است که بر اساس آن، هنگامی که یک گروه
	تحت ستم سیستماتیک قرار دارد و همه‌ی راه‌های داخلی برای احقاق حق
	بسته شده‌اند، آن گروه حق جدایی از واحد سیاسی موجود و تشکیل
	واحد مستقل را به دست
	می‌آورد\src{\lr{Buchanan, 2004, pp.\,331--400}}.
\end{tcolorbox}

\thinker{آلن بوکانان}%
\footnote{%
	آلن بوکانان (\lr{Allen Buchanan, 1948–})؛%
	فیلسوف سیاسی آمریکایی و نظریه‌پرداز جدایی‌طلبی.%
}%
پنج شرط برای مشروعیت جدایی جبرانی پیشنهاد کرده
است\src{\lr{Buchanan, 2004, pp.\,331--400}}:

\begin{enumerate}[label=\textbf{\arabic*.}, nosep, rightmargin=1em]
	\item ستم سیستماتیک و مستمر علیه گروه؛
	\item ناممکنی اصلاح از درون (شکست گفت‌وگو و اصلاحات)؛
	\item وجود سرزمین مشخص و جمعیت منسجم؛
	\item توانایی تشکیل دولت کارآمد؛
	\item عدم ایجاد ستم جدید علیه اقلیت‌های درون گروه جدید.
\end{enumerate}

\subsection{زنجیره‌ی خشونت: وقتی احقاق حق خود منبع خشونت می‌شود}
\label{subsec:ch02-t4-violence}

\begin{tcolorbox}[enemybox,
	title={معضل زنجیره‌ی خشونت}]
	مشکل آن است که فرایند جدایی جبرانی، حتی هنگامی که از نظر
	اخلاقی و حقوقی موجه باشد، اغلب \textbf{زنجیره‌ای از خشونت}
	را آغاز می‌کند:
	\begin{enumerate}[label=\textbf{\alph*.}, nosep, rightmargin=1em]
		\item دولت مرکزی با خشونت واکنش نشان می‌دهد؛
		\item گروه جدایی‌طلب نیز مسلح می‌شود؛
		\item نیروهای خارجی وارد می‌شوند؛
		\item غیرنظامیان بیشترین آسیب را می‌بینند؛
		\item پس از جدایی، اقلیت‌های درون واحد جدید خود سرکوب می‌شوند.
	\end{enumerate}
	
	نمونه‌ها: یوگسلاوی (۱۹۹۱--۱۹۹۹)، سودان جنوبی (۲۰۱۱ تا کنون)،
	اقلیم کردستان عراق (دهه‌ی ۱۹۹۰ تا کنون).
\end{tcolorbox}

\subsubsection{تقصیر، فداکاری و هزینه از جیب دیگری}
\label{subsubsec:ch02-t4-blame}

یکی از پیچیده‌ترین ابعاد اخلاقی این وضعیت آن است که \textbf{تصمیم
	به مقاومت یا جدایی اغلب توسط نخبگان گرفته می‌شود، اما هزینه‌اش
	را مردم عادی می‌پردازند}.
\thinker{جف مک‌ماهان}%
\footnote{%
	جف مک‌ماهان (\lr{Jeff McMahan, 1954–})؛%
	فیلسوف اخلاق جنگ در دانشگاه آکسفورد.%
}%
استدلال می‌کند که «حق به جنگ» (\lr{jus ad bellum}) باید
با «عدالت کسانی که هزینه می‌پردازند» سنجیده
شود\src{\lr{McMahan, 2009, pp.\,38--65}}.

\begin{tcolorbox}[wavebox, title={پرسش اخلاقی بنیادین}]
	آیا رهبران یک جنبش حق دارند مردم خود را به جنگی بکشانند
	که هزینه‌ی آن عمدتاً بر دوش غیرنظامیان است؟ و آیا بازیگر
	خارجی که با حمایت از جنبش، آتش جنگ را شعله‌ور می‌کند، در
	قبال این هزینه‌ها مسئول نیست؟
	
	این مسئله ما را به مفهوم \concept{تقصیر توزیعی}%
	\footnote{تقصیر توزیعی؛ به انگلیسی: \lr{Distributed Culpability}}%
	می‌رساند: مسئولیت رنج ناشی از مداخله یا جدایی میان دولت
	سرکوبگر، جنبش جدایی‌طلب، نیروی خارجی مداخله‌گر و حتی جامعه‌ی
	بین‌المللی ناظر توزیع
	می‌شود\src{\lr{Raz, 1986, pp.\,369--395}}.
\end{tcolorbox}


% ════════════════════════════════════════════════════════════════
\section{تنش پنجم: گروه پیشرو، بازیگر خارجی و مشروعیت مداخله}
\label{sec:ch02-tension5}

\subsection{ساختار مسئله}
\label{subsec:ch02-t5-structure}

فرض کنید در جامعه‌ای، \concept{گروه پیشرو}%
\footnote{منظور گروهی است که ارزش‌های غیرسنتی و نوآورانه، مانند
	حقوق زنان، سکولاریسم، آزادی بیان را تعقیب می‌کند.}%
وجود دارد که ارزش‌هایش با یک قدرت خارجی همسوست. آن قدرت خارجی
ممکن است:

\begin{enumerate}[label=\textbf{\alph*.}, nosep, rightmargin=1em]
	\item واقعاً به این ارزش‌ها متعهد باشد (ایده‌آلیسم)؛
	\item این ارزش‌ها بخشی از پروپاگاندا و تفوق‌طلبی‌اش باشد
	(امپریالیسم فرهنگی)؛
	\item ترکیبی از هر دو (واقع‌گرایی ایده‌آلیستی).
\end{enumerate}

هنگامی که این قدرت خارجی «به دعوت» گروه پیشرو به گروه سنتی حمله
می‌کند، چند پرسش مطرح می‌شود:

\begin{tcolorbox}[mybox, title={پرسش‌های کلیدی تنش پنجم}]
	\begin{enumerate}[label=\textbf{\arabic*.}, nosep, rightmargin=1em]
		\item آیا گروه پیشرو نماینده‌ی جامعه است یا اقلیت نخبگانی؟
		\item آیا ارزش‌های مورد ادعای بازیگر خارجی واقعی‌اند یا ابزاری؟
		\item حتی اگر ارزش‌ها واقعی باشند، آیا تحمیل نظامی آن‌ها
		مشروع است؟
		\item آیا طرفداری از یک طرف به سرکوب طرف مقابل ختم نمی‌شود
		(معکوس‌سازی ستم)؟
		\item آیا عنصر مشترکی وجود دارد که هم تکثر فرهنگی را حفظ کند
		و هم حد و حدودی برای اجبار بشناسد؟
	\end{enumerate}
\end{tcolorbox}

\subsection{خطر معکوس‌سازی ستم}
\label{subsec:ch02-t5-reversal}

\thinker{رینهولد نیبور}%
\footnote{%
	رینهولد نیبور (\lr{Reinhold Niebuhr, 1892--1971})؛%
	الهی‌دان و فیلسوف سیاسی آمریکایی.%
}%
هشدار داده که «کسانی که ستم می‌بینند، هنگامی که به قدرت برسند،
اغلب ستم
می‌کنند»\src{\lr{Niebuhr, 1932, pp.\,xi--xii}}.
تاریخ مداخلات خارجی مملو از نمونه‌هایی است که در آن‌ها «آزادسازی»
صرفاً به \concept{جابه‌جایی ستمگر} انجامیده: ستمگر قدیمی رفته و
ستمگر جدید، اغلب وابسته به قدرت مداخله‌گر، جایش را گرفته است.

\subsection{جستجوی عنصر مشترک: حداقل‌های عدالت}
\label{subsec:ch02-t5-minimum}

\thinker{جان رالز}%
\footnote{جان رالز (\lr{John Rawls, 1921--2002}).}%
در کتاب قانون مردمان مفهوم
\concept{حداقل‌های عدالت}%
\footnote{حداقل‌های عدالت؛ به انگلیسی: \lr{Decent Minimum}}%
را مطرح می‌کند: مجموعه‌ای از حقوق بنیادین که هیچ جامعه‌ای،
فارغ از فرهنگ و سنتش، نباید از آن‌ها عدول
کند\src{\lr{Rawls, 1999, pp.\,64--78}}.
این حداقل‌ها شامل حق حیات، آزادی از بردگی، آزادی از شکنجه و
حداقلی از آزادی وجدان هستند.

\begin{tcolorbox}[wavebox, title={پیشنهاد تحلیلی}]
	عنصر مشترکی که هم تکثر فرهنگی را حفظ کند و هم اجبار را محدود
	سازد، می‌تواند بر \textbf{سه اصل} استوار باشد:
	\begin{enumerate}[label=\textbf{\arabic*.}, nosep, rightmargin=1em]
		\item \concept{حق خروج}%
		\footnote{حق خروج؛ به انگلیسی: \lr{Right of Exit}}: هر فرد
		باید حق واقعی (نه صرفاً نظری) ترک گروه یا فرهنگ خود را
		داشته باشد. اگر فرهنگی اجازه‌ی خروج نمی‌دهد، آن فرهنگ نه
		یک «انتخاب آزاد» بلکه یک «زندان هویتی»
		است\src{\lr{Kukathas, 2003, pp.\,96--133}}؛
		\item \concept{اصل عدم سلطه}%
		\footnote{عدم سلطه؛ به انگلیسی: \lr{Non-Domination}}:
		هیچ گروهی، نه سنتی و نه مدرن، نباید بتواند اراده‌ی خود
		را بر دیگران به‌گونه‌ای خودسرانه تحمیل کند. این اصل را
		\thinker{فیلیپ پتیت}%
		\footnote{فیلیپ پتیت (\lr{Philip Pettit, 1945–})؛
			فیلسوف سیاسی ایرلندی-استرالیایی.}%
		مبنای آزادی جمهوریخواهانه قرار داده
		است\src{\lr{Pettit, 1997, pp.\,51--79}}؛
		\item \concept{اصل تناسب}%
		\footnote{تناسب؛ به انگلیسی: \lr{Proportionality}}:
		هرگونه مداخله (چه داخلی و چه خارجی) باید متناسب با شدت ستم
		باشد. واکنش نظامی تنها در صورت ناکارآمدی همه‌ی ابزارهای
		غیرنظامی مشروع
		است\src{\lr{Hurka, 2005, pp.\,34--54}}.
	\end{enumerate}
\end{tcolorbox}


% ════════════════════════════════════════════════════════════════
\section{تنش ششم: بلوغ درون‌زاد ارزش‌ها و ناممکنی تزریق}
\label{sec:ch02-tension6}

\subsection{ارزش‌هایی که باید رُشد کنند، نه کاشته شوند}
\label{subsec:ch02-t6-growth}

\begin{tcolorbox}[quotebox, title={}]
	«دموکراسی را نمی‌توان صادر کرد. دموکراسی باید از دل تجربه‌ی
	یک ملت بروید، با همه‌ی دردها و خطاهایش.»
	
	\hfill ---واتسلاو هاول\src{\lr{Havel, 1991, p.\,18}}
\end{tcolorbox}

شاید مهم‌ترین اعتراض علیه مداخلات «ارزش‌محور» آن باشد که
ارزش‌هایی مانند
\concept{رواداری}%
\footnote{رواداری؛ به انگلیسی: \lr{Tolerance}}،
\concept{تکثرگرایی}%
\footnote{تکثرگرایی؛ به انگلیسی: \lr{Pluralism}}،
\concept{خودبنیادی}%
\footnote{خودبنیادی؛ به انگلیسی: \lr{Autonomy / Self-Foundation}}
و \concept{پذیرش خطاپذیری خود}%
\footnote{خطاپذیری؛ به انگلیسی: \lr{Fallibilism}}
ماهیتی دارند که با \textbf{تزریق} ناسازگار است.

\thinker{یورگن هابرماس}%
\footnote{%
	یورگن هابرماس (\lr{Jürgen Habermas, 1929–})؛%
	فیلسوف و جامعه‌شناس آلمانی.%
}%
استدلال می‌کند که نهادهای دموکراتیک تنها هنگامی پایدارند که بر
\concept{فرهنگ سیاسی}%
\footnote{فرهنگ سیاسی؛ به انگلیسی: \lr{Political Culture}}%
متناسبی استوار باشند؛ فرهنگی که از طریق تجربه‌ی عملی گفت‌وگو،
تضاد و سازش شکل
می‌گیرد\src{\lr{Habermas, 1996, pp.\,287--328}}.

\subsection{شواهد تاریخی: تفاوت نهادسازی درون‌زاد و برون‌زاد}
\label{subsec:ch02-t6-evidence}

\begin{table}[htbp]
	\centering
	\caption{مقایسه‌ی نهادسازی درون‌زاد و برون‌زاد دموکراتیک}
	\label{tab:ch02-endogenous}
	\begin{adjustbox}{max width=\textwidth}
		\begin{tabular}{>{\bfseries\color{PrimaryMid}}p{2.2cm}
				p{2.8cm} p{2.5cm} p{2cm} p{3cm}}
			\toprule
			\textbf{نمونه} &
			\textbf{نوع فرایند} &
			\textbf{مدت بلوغ} &
			\textbf{پایداری} &
			\textbf{بحران‌ها} \\
			\midrule
			انگلستان &
			عمدتاً درون‌زاد &
			چند سده &
			بسیار بالا &
			جنگ داخلی، ولی عبور موفق \\
			\midrule
			فرانسه &
			درون‌زاد ولی خشونت‌بار &
			بیش از یک سده &
			بالا (پس از تثبیت) &
			چندین انقلاب و بازگشت \\
			\midrule
			ژاپن &
			برون‌زاد (اشغال آمریکا) &
			نسبتاً سریع &
			بالا &
			تنش‌های هویتی پایدار \\
			\midrule
			عراق &
			برون‌زاد (اشغال آمریکا) &
			ناتمام &
			بسیار پایین &
			جنگ داخلی، داعش \\
			\midrule
			افغانستان &
			برون‌زاد (اشغال ناتو) &
			شکست &
			صفر &
			بازگشت طالبان \\
			\bottomrule
		\end{tabular}
	\end{adjustbox}
\end{table}

تفاوت ژاپن و عراق/افغانستان نشان می‌دهد که نهادسازی برون‌زاد
تنها در شرایط بسیار خاص موفق است.
\thinker{فرانسیس فوکویاما}%
\footnote{فرانسیس فوکویاما (\lr{Francis Fukuyama, 1952–}).}%
در کتاب دولت‌سازی (\lr{\textit{State-Building}}، ۲۰۰۴) استدلال
می‌کند که دموکراسی‌سازی از بیرون نیازمند وجود پیشینی
\concept{ظرفیت دولتی}%
\footnote{ظرفیت دولتی؛ به انگلیسی: \lr{State Capacity}}
و \concept{سرمایه‌ی اجتماعی}%
\footnote{سرمایه‌ی اجتماعی؛ به انگلیسی: \lr{Social Capital}}
در جامعه‌ی هدف
است\src{\lr{Fukuyama, 2004, pp.\,40--58}}.

\subsection{فرایند درونی‌سازی: زمان، تضاد و عبور}
\label{subsec:ch02-t6-internalization}

\begin{tcolorbox}[mybox,
	title={مراحل درونی‌سازی ارزش‌های دموکراتیک}]
	بر اساس تلفیقی از نظریات
	\thinker{هابرماس}، \thinker{رالز}
	و \thinker{چارلز تیلی}%
	\footnote{چارلز تیلی (\lr{Charles Tilly, 1929--2008})؛
		جامعه‌شناس تاریخی آمریکایی.}،
	فرایند درونی‌سازی ارزش‌ها عموماً شامل مراحل زیر است:
	
	\begin{enumerate}[label=\textbf{مرحله‌ی \arabic*:}, nosep, rightmargin=1em]
		\item \concept{مواجهه}: آشنایی اولیه با ایده‌های جدید (اغلب
		از طریق تماس با بیرون یا نخبگان آشنا به بیرون)؛
		\item \concept{تضاد}: برخورد ایده‌های جدید با ساختارهای موجود
		و مقاومت نهادی-فرهنگی؛
		\item \concept{تجربه‌ی عملی}: آزمون‌وخطا در پیاده‌سازی نهادهای
		جدید (اغلب همراه با شکست‌های مکرر)؛
		\item \concept{سازش و تعادل}: یافتن ترکیبی بومی از ارزش‌های
		قدیم و جدید که عملاً کار کند؛
		\item \concept{نهادینه‌شدن}: تبدیل ارزش‌های جدید به «بدیهیات»
		فرهنگ سیاسی.
	\end{enumerate}
	
	\textbf{این فرایند راه میانبری ندارد.} هرگونه تلاش برای جهش از
	مرحله‌ی ۱ به مرحله‌ی ۵، چه از درون و چه از بیرون،
	به \concept{نهادهای صوری و شکننده}
	می‌انجامد\src{\lr{Tilly, 2007, pp.\,1--24; Huntington, 1968, pp.\,1--8}}.
\end{tcolorbox}

\begin{figure}[htbp]
	\centering
	\begin{tikzpicture}[
		>=Stealth,
		every node/.style={font=\small},
		stage/.style={
			draw=PrimaryMid, fill=BgCool, rounded corners=4pt,
			minimum width=2.4cm, minimum height=1.4cm,
			text width=2.2cm, align=center, font=\small\bfseries
		}
		]
		% مراحل
		\node[stage] (s1) at (0,0)   {\rl{مواجهه}};
		\node[stage] (s2) at (3.5,0) {\rl{تضاد}};
		\node[stage] (s3) at (7,0)   {\rl{تجربه‌ی}\\[1pt]{\footnotesize\rl{عملی}}};
		\node[stage] (s4) at (10.5,0){\rl{سازش و}\\[1pt]{\footnotesize\rl{تعادل}}};
		\node[stage] (s5) at (14,0)  {\rl{نهادینه‌}\\[1pt]{\footnotesize\rl{شدن}}};
		
		% فلش‌های اصلی
		\draw[->, thick, AccentGold] (s1) -- (s2);
		\draw[->, thick, AccentGold] (s2) -- (s3);
		\draw[->, thick, AccentGold] (s3) -- (s4);
		\draw[->, thick, AccentGold] (s4) -- (s5);
		
		% فلش بازگشت (شکست)
		\draw[->, thick, AccentRed, dashed, bend left=40]
		(s3.north) to node[above, font=\scriptsize, text=AccentRed]
		{\rl{شکست و بازگشت}} (s2.north);
		
		% فلش جهش ناموفق
		\draw[->, thick, AccentRed, dotted, bend right=30]
		(s1.south) to node[below, font=\scriptsize, text=AccentRed,
		text width=3cm, align=center]
		{\rl{جهش تحمیلی}\\[1pt]{\scriptsize\rl{(صوری و شکننده)}}} (s5.south);
		
		% زمان
		\draw[->, very thick, NeutralDark]
		(-1.5,-2.2) -- (15.5,-2.2);
		\node[font=\footnotesize] at (14.5,-2.5) {\rl{زمان (نسل‌ها)}};
		
		% عنوان
		\node[above, font=\normalsize\bfseries, text=PrimaryDeep]
		at (7,2) {\rl{فرایند درونی‌سازی ارزش‌های دموکراتیک}};
	\end{tikzpicture}
	\caption{فرایند درونی‌سازی ارزش‌ها: مسیر طبیعی و جهش تحمیلی}
	\label{fig:ch02-internalization}
\end{figure}


% ════════════════════════════════════════════════════════════════
\section{تنش هفتم: جامعه‌ی دوپاره و شکنندگی دموکراسی}
\label{sec:ch02-tension7}

\subsection{قطبی‌شدن و بن‌بست ساختاری}
\label{subsec:ch02-t7-polarization}

\begin{tcolorbox}[enemybox,
	title={معضل جوامع عمیقاً دوپاره}]
	هنگامی که یک جامعه به دو (یا چند) اردوگاه بنیادین تقسیم شده،
	مثلاً \concept{سنتی/مدرن}، \concept{سکولار/مذهبی}،
	\concept{توسعه‌گرا/هویت‌گرا}، دموکراسی به معنای متعارف
	(حکومت اکثریت) نه تنها راه‌حل نیست، بلکه ممکن است \textbf{خود
		بخشی از مشکل} باشد.
\end{tcolorbox}

% ═══════════════════════════════════════════════════════════════
%  ch02.tex — ادامه (از تنش هفتم: لیپهارت)
% ═══════════════════════════════════════════════════════════════

\thinker{آرنت لیپهارت}%
\footnote{%
	آرنت لیپهارت (\lr{Arend Lijphart, 1936–})؛%
	دانشمند علوم سیاسی هلندی-آمریکایی.%
}%
با مطالعه‌ی تطبیقی دموکراسی‌های چندپاره (هلند، بلژیک، لبنان،
سوئیس) نشان داده که در جوامع
\concept{چنددسته‌ای}%
\footnote{جوامع چنددسته‌ای؛ به انگلیسی:
	\lr{Consociational / Segmented Societies}}%
دموکراسی اکثریتی (\lr{Majoritarian}) به بی‌ثباتی و سرکوب
اقلیت‌ها می‌انجامد. او بدیل
\concept{دموکراسی توافقی}%
\footnote{دموکراسی توافقی؛ به انگلیسی:
	\lr{Consociational Democracy}}%
را پیشنهاد می‌کند که بر اصول وتوی متقابل، ائتلاف بزرگ،
خودمختاری بخشی و نمایندگی تناسبی استوار
است\src{\lr{Lijphart, 1977, pp.\,25--52}}.

\subsection{پوپولیسم: وقتی دموکراسی علیه خود عمل می‌کند}
\label{subsec:ch02-t7-populism}

\begin{tcolorbox}[defbox,
	title={مفهوم: پوپولیسم در جوامع دوپاره}]
	\concept{پوپولیسم}%
	\footnote{پوپولیسم؛ به انگلیسی: \lr{Populism}}%
	در جوامع عمیقاً قطبی‌شده ویژگی خاصی می‌یابد: رهبر پوپولیست
	با ادعای نمایندگی «مردم واقعی» در برابر «نخبگان فاسد» یا
	«دشمنان ملت»، شکاف موجود را عمیق‌تر می‌کند و دموکراسی را
	از درون تهی
	می‌سازد\src{\lr{Müller, 2016, pp.\,1--6; Mudde, 2004, pp.\,541--563}}.
\end{tcolorbox}

\thinker{یان-ورنر مولر}%
\footnote{%
	یان-ورنر مولر (\lr{Jan-Werner Müller, 1970–})؛%
	نظریه‌پرداز سیاسی آلمانی در دانشگاه پرینستون.%
}%
استدلال می‌کند که پوپولیسم ذاتاً \concept{ضدتکثرگرایانه} است:
پوپولیست مدعی است «فقط من» نماینده‌ی «مردم واقعی» هستم و هر
مخالفتی غیرمشروع
است\src{\lr{Müller, 2016, pp.\,19--42}}.
در جامعه‌ی دوپاره، این ادعا به‌سادگی با یکی از دو قطب
همسو می‌شود و نیمه‌ی دیگر را حذف می‌کند.

\subsection{مداخله‌ی خارجی در جوامع دوپاره: تشدید یا تعدیل؟}
\label{subsec:ch02-t7-intervention}

تجربه‌ی تاریخی نشان می‌دهد که مداخله‌ی خارجی در جوامع دوپاره
اغلب \textbf{تشدیدکننده} بوده، نه تعدیل‌کننده:

\begin{table}[htbp]
	\centering
	\caption{تأثیر مداخله‌ی خارجی بر جوامع دوپاره}
	\label{tab:ch02-polarized}
	\begin{adjustbox}{max width=\textwidth}
		\begin{tabular}{>{\bfseries\color{PrimaryMid}}p{2.5cm}
				p{3cm} p{3cm} p{4cm}}
			\toprule
			\textbf{نمونه} & \textbf{شکاف اصلی} &
			\textbf{مداخله‌گر} & \textbf{نتیجه} \\
			\midrule
			لبنان ۱۹۷۵--۱۹۹۰ &
			مسیحی/مسلمان &
			سوریه، اسرائیل، ایران &
			تعمیق شکاف، جنگ ۱۵ ساله \\
			\midrule
			عراق ۲۰۰۳--اکنون &
			شیعه/سنی/کرد &
			آمریکا، ایران &
			جنگ فرقه‌ای، ظهور داعش \\
			\midrule
			اوکراین ۲۰۱۴--اکنون &
			غربگرا/روسگرا &
			روسیه، ناتو &
			تجزیه، جنگ تمام‌عیار \\
			\midrule
			لیبی ۲۰۱۱--اکنون &
			قبایل/مناطق &
			ناتو، ترکیه، مصر، امارات &
			دولت ناکام، دوگانگی قدرت \\
			\bottomrule
		\end{tabular}
	\end{adjustbox}
\end{table}

\begin{tcolorbox}[wavebox,
	title={استثناء: بوسنی و مکانیسم دیتون}]
	توافقنامه‌ی دیتون (۱۹۹۵) نمونه‌ای نادر از مداخله‌ی خارجی در
	جامعه‌ی چندپاره است که، علیرغم نقص‌های فراوان، به خاتمه‌ی
	جنگ و ایجاد ساختار سیاسی نسبتاً پایدار انجامید. رمز نسبی موفقیت
	آن، پذیرش واقعیت چندپارگی و طراحی نهادهای \concept{توافقی}
	به‌جای تحمیل دموکراسی اکثریتی
	بود\src{\lr{Bieber, 2006, pp.\,1--28}}.
	اما حتی بوسنی نیز پس از سه دهه هنوز جامعه‌ای عمیقاً تقسیم‌شده
	باقی مانده است.
\end{tcolorbox}


% ════════════════════════════════════════════════════════════════
\section{تنش هشتم: بازی‌ای که نمی‌توان متوقف کرد}
\label{sec:ch02-tension8}

\subsection{«گر رسم شود که مست گیرند، در شهر هر آنکه هست گیرند»}
\label{subsec:ch02-t8-proverb}

\begin{tcolorbox}[quotebox, title={}]
	{\large
		گر رسم شود که مست گیرند
		
		در شهر هر آنکه هست گیرند}
	
	\hfill ---منسوب به حافظ
\end{tcolorbox}

این بیت ناظر بر یک حقیقت عمیق سیاسی است: \textbf{هنگامی که
	مشروعیت یک اقدام خارق‌العاده (مداخله، انقلاب، سرکوب) پذیرفته
	شود، محدود کردن دامنه‌ی آن تقریباً ناممکن می‌شود}. اگر مداخله‌ی
خارجی «برای نجات» مشروع شناخته شود، چه چیزی مانع از آن است که
هر قدرتی «نجات» را بهانه‌ی مداخله در هر جایی قرار دهد؟

\subsection{مسئله‌ی سابقه و تسری}
\label{subsec:ch02-t8-precedent}

\thinker{فرناندو تسون}%
\footnote{%
	فرناندو تسون (\lr{Fernando R.\,Tesón, 1950–})؛%
	حقوقدان بین‌المللی آرژانتینی.%
}%
که خود از مدافعان سرسخت مداخله‌ی بشردوستانه است، اذعان
می‌کند که مهم‌ترین اعتراض علیه حق مداخله،
\concept{مسئله‌ی سابقه}%
\footnote{مسئله‌ی سابقه؛ به انگلیسی: \lr{Precedent Problem}}%
است: هرگونه مشروعیت‌بخشی به مداخله، سابقه‌ای می‌سازد که
قدرت‌های فرصت‌طلب از آن سوءاستفاده خواهند
کرد\src{\lr{Tesón, 2005, pp.\,153--174}}.

نمونه‌های تاریخی:
\begin{enumerate}[label=\textbf{\arabic*.}, nosep, rightmargin=1em]
	\item مداخله‌ی ناتو در کوزوو (۱۹۹۹) به‌عنوان سابقه‌ای برای
	توجیه حمله‌ی روسیه به گرجستان (۲۰۰۸) و اوکراین (۲۰۱۴)؛
	\item مفهوم \lr{R2P}%
	\footnote{مسئولیت حمایت؛ به انگلیسی:
		\lr{Responsibility to Protect}.}
	در لیبی (۲۰۱۱) به‌عنوان سابقه‌ای که چین و روسیه سپس مانع
	تکرار آن در سوریه شدند؛
	\item «مداخله‌ی بشردوستانه» هند در بنگلادش (۱۹۷۱) به‌عنوان
	سابقه‌ای برای مداخلات بعدی در سریلانکا.
\end{enumerate}

\subsection{صورت‌بندی دقیق: شرایط لازم برای محدودسازی}
\label{subsec:ch02-t8-formulation}

\begin{tcolorbox}[legalbox,
	title={شرایط محدودسازی حق مداخله: پیشنهاد ترکیبی}]
	بر اساس ترکیب آرای \thinker{تسون}، \thinker{والزر} و
	گزارش کمیسیون بین‌المللی مداخله و حاکمیت ملی%
	\footnote{کمیسیون بین‌المللی مداخله و حاکمیت ملی؛ به انگلیسی:
		\lr{ICISS}, ۲۰۰۱.}،
	شرایط زیر برای مشروعیت مداخله‌ی بشردوستانه پیشنهاد
	می‌شود\src{\lr{ICISS, 2001, pp.\,32--37}}:
	
	\begin{enumerate}[label=\textbf{\arabic*.}, nosep, rightmargin=1em]
		\item \concept{آستانه‌ی ستم}: فقط نسل‌کشی، پاکسازی قومی،
		جنایات جنگی و جنایات علیه بشریت (نه هر نقض حقوق بشر)؛
		\item \concept{نیت درست}: هدف اصلی (نه لزوماً تنها) باید توقف
		رنج انسانی باشد؛
		\item \concept{آخرین چاره}: همه‌ی ابزارهای دیپلماتیک و اقتصادی
		امتحان شده یا آشکارا ناکارآمد باشند؛
		\item \concept{تناسب ابزار}: ابزار نظامی باید حداقلی و متناسب
		با هدف باشد؛
		\item \concept{احتمال معقول موفقیت}: مداخله نباید وضع را بدتر
		کند؛
		\item \concept{مرجع مشروع}: ترجیحاً شورای امنیت سازمان ملل
		یا ائتلاف منطقه‌ای گسترده.
	\end{enumerate}
\end{tcolorbox}

\subsection{نقد: آیا صورت‌بندی دقیق ممکن است؟}
\label{subsec:ch02-t8-critique}

\thinker{نوام چامسکی}%
\footnote{نوام چامسکی (\lr{Noam Chomsky, 1928–}).}%
به‌تندی استدلال می‌کند که هیچ صورت‌بندی‌ای نمی‌تواند از سوءاستفاده
جلوگیری کند، زیرا قدرت‌های بزرگ همواره قواعد را به نفع خود
تفسیر می‌کنند\src{\lr{Chomsky, 1999, pp.\,34--67}}.
از سوی دیگر،
\thinker{گرت اونز}%
\footnote{%
	گرت اونز (\lr{Gareth Evans, 1944–})؛%
	وزیر خارجه‌ی سابق استرالیا و از معماران \lr{R2P}.%
}%
معتقد است که «وجود خطر سوءاستفاده نمی‌تواند توجیه‌گر بی‌عملی
در برابر نسل‌کشی باشد»\src{\lr{Evans, 2008, pp.\,31--55}}.

\begin{tcolorbox}[mybox,
	title={نتیجه‌ی تحلیلی: پارادوکس غیرقابل حل}]
	حقیقت آن است که تنش هشتم، بازی‌ای که نمی‌توان متوقف کرد،
	\textbf{پارادوکسی بنیاداً غیرقابل حل} است. هرگونه مشروعیت‌بخشی
	به مداخله، خطر تسری و سوءاستفاده دارد؛ و هرگونه ممنوعیت مطلق
	مداخله، خطر تداوم نسل‌کشی و ستم سیستماتیک را دارد. آنچه می‌توان
	کرد، نه «حل» مسئله بلکه \concept{مدیریت آگاهانه‌ی آن} است:
	\begin{itemize}[nosep, rightmargin=1em]
		\item تدوین معیارهای دقیق و شفاف (هرچند همواره ناکامل)؛
		\item ایجاد نهادهای بین‌المللی پاسخگو (هرچند همواره ناقص)؛
		\item ارتقای هزینه‌ی سیاسی سوءاستفاده از حق مداخله؛
		\item تقویت مکانیسم‌های درون‌زاد حل تضاد.
	\end{itemize}
\end{tcolorbox}


% ════════════════════════════════════════════════════════════════
\section{ترکیب تنش‌ها: ماتریس تعامل}
\label{sec:ch02-matrix}

هشت تنشی که بررسی شد، مستقل از یکدیگر نیستند. آن‌ها در
وضعیت‌های واقعی همزمان عمل می‌کنند و ترکیب‌های مختلفی از معضلات
را تولید می‌کنند.

\begin{landscape}
	\begin{table}[htbp]
		\centering
		\caption{ماتریس تعامل هشت تنش بنیادین نجات‌بخشی/نجات‌طلبی}
		\label{tab:ch02-tension-matrix}
		\begin{adjustbox}{max width=\linewidth}
			\begin{tabular}{>{\bfseries\scriptsize\color{PrimaryMid}}p{1.3cm}
					>{\scriptsize}p{2cm}
					>{\scriptsize}p{2cm}
					>{\scriptsize}p{2cm}
					>{\scriptsize}p{2cm}
					>{\scriptsize}p{2cm}
					>{\scriptsize}p{2cm}
					>{\scriptsize}p{2cm}
					>{\scriptsize}p{2cm}}
				\toprule
				& \textbf{ت۱} & \textbf{ت۲}
				& \textbf{ت۳} & \textbf{ت۴}
				& \textbf{ت۵} & \textbf{ت۶}
				& \textbf{ت۷} & \textbf{ت۸} \\
				\midrule
				ت۱ & --- &
				اجبار ساختاری مشروع؟ &
				کدام ساختار؟ &
				ساختارشکنی &
				چه کسی می‌سازد؟ &
				بدون بلوغ پایدار نیست &
				برای کدام نیمه؟ &
				سابقه‌سازی \\
				\midrule
				ت۲ & &  --- &
				اجبار به مدرنیته &
				اجبار و خشونت &
				بازیگر خارجی اجبار می‌کند &
				اجبار مانع بلوغ &
				اجبار بر اقلیت &
				تعمیم اجبار \\
				\midrule
				ت۳ & & & --- &
				جدایی فرهنگی &
				مدرنیته‌ی تحمیلی &
				مدرنیته‌ی درون‌زاد؟ &
				شکاف سنت/مدرنیته &
				تسری مدل غربی \\
				\midrule
				ت۴ & & & & --- &
				مداخله در جنگ داخلی &
				خشونت مانع بلوغ &
				تعمیق دوپارگی &
				سابقه‌ی جدایی \\
				\midrule
				ت۵ & & & & & --- &
				پیشرو: عامل یا مانع؟ &
				پیشرو = یک قطب &
				تسری حمایت \\
				\midrule
				ت۶ & & & & & & --- &
				بلوغ نامتقارن &
				تسری نسخه‌ی واحد \\
				\midrule
				ت۷ & & & & & & & --- &
				سابقه‌ی مدیریت شکاف \\
				\midrule
				ت۸ & & & & & & & & --- \\
				\bottomrule
			\end{tabular}
		\end{adjustbox}
	\end{table}
\end{landscape}


% ════════════════════════════════════════════════════════════════
\section{جمع‌بندی: نه نجات‌بخشی آسان، نه نجات‌طلبی بی‌هزینه}
\label{sec:ch02-conclusion}

تحلیل هشت تنش بنیادین نشان می‌دهد که مسئله‌ی نجات‌بخشی و نجات‌طلبی
ساده‌انگاری برنمی‌تابد. نه آرمان‌گرایی ناب (مداخله‌ی خارجی
همیشه خوب است) و نه بدبینی مطلق (مداخله‌ی خارجی همیشه بد است)
با واقعیت سازگارند. آنچه این بررسی آشکار می‌کند، عبارت است از:

\begin{enumerate}[label=\textbf{\arabic*.}, nosep, rightmargin=1em]
	\item رهایی‌بخشی ساختاری از بیرون \textbf{ممکن} اما
	\textbf{مشروط} است (تنش اول)؛
	\item اجبار به آزادی \textbf{تنها در حداقل‌ها و به‌صورت موقت}
	قابل دفاع است (تنش دوم)؛
	\item تحمیل مدرنیته‌ی خطی \textbf{نه ممکن و نه مطلوب} است؛
	اما «ترجیحات تطبیقی» نیز نباید بی‌نقد پذیرفته شوند
	(تنش سوم)؛
	\item حق جدایی جبرانی \textbf{وجود دارد} اما
	\textbf{زنجیره‌ی خشونت} آن ناگزیر است و «تقصیر» توزیعی
	است (تنش چهارم)؛
	\item حمایت از گروه پیشرو \textbf{ممکن} است اما نباید به
	\textbf{سرکوب معکوس} بینجامد (تنش پنجم)؛
	\item ارزش‌های دموکراتیک \textbf{باید درون‌زاد شوند}؛ تزریق
	بیرونی فقط نهادهای صوری تولید می‌کند (تنش ششم)؛
	\item در جوامع دوپاره، \textbf{دموکراسی اکثریتی کافی نیست}
	و مکانیسم‌های توافقی ضرورت دارند (تنش هفتم)؛
	\item مشروعیت‌بخشی به مداخله \textbf{سابقه‌ساز} است و
	صورت‌بندی دقیق شرایط آن ضروری ولی همواره ناقص
	(تنش هشتم).
\end{enumerate}

در فصل‌های آتی، این چارچوب تحلیلی هشت‌گانه بر مطالعات موردی
تاریخی اعمال خواهد شد؛ ابتدا بر تاریخ ایران و سپس بر
تجربیات جهانی. هدف آن است که ببینیم این تنش‌ها در عمل چگونه
ظاهر شده و جوامع مختلف چگونه با آن‌ها مواجه شده‌اند.